% ===========================================================================
% SOURCES: docs/claude_project/03_HARDWARE_ARCHITECTURE.md
%          docs/controller_design.md (Greek - the architecture study)
%          docs/pointwise_datapath_explained.md
%          docs/latency_map.md
%          docs/design_decisions.md: DD-013..DD-018, DD-021, DD-029
%          hardware/rtl/**
%
% This chapter describes WHAT WAS BUILT and why. It does NOT contain the
% timing-closure narrative (Chapter~\ref{chap:timing}) or the verification method
% (Chapter~\ref{chap:verification}) - both were split out so this chapter stays
% readable. Resist pulling them back in.
% ===========================================================================
\chapter{Accelerator Architecture}
\label{chap:architecture}

This chapter presents the architecture of the proposed accelerator and the
design decisions behind it. Section~\ref{sec:arch-overview} first gives an
overview of the complete system and summarizes its dataflow, and
Section~\ref{sec:engines-philosophy} states the separation of concerns that
shapes every engine described after it. The remaining sections justify the
chosen degree of parallelism (Section~\ref{sec:parallelism}), the on-chip
memory organization (Section~\ref{sec:memory}), the control unit and its
instruction set (Section~\ref{sec:control}), and the latency contract that
binds the feeders to the engines (Section~\ref{sec:latency-contract}). The
integration of the accelerator into the Zynq UltraScale+ MPSoC is deferred to
Chapter~\ref{chap:integration}.

\section{Proposed Accelerator Architecture}
\label{sec:arch-overview}
% The block diagram, and the one-paragraph summary of the dataflow. A reader
% should be able to stop after this section and still know what the thing is.
% The figure below is deliberately the SIMPLE one: engines named by what they
% compute, one shared output stage, one pool, and the loop. Every width, every
% tiling factor and the accel_top/accel_sys/accel_axi layering are held back for
% the sections that argue them (\ref{sec:parallelism}, \ref{sec:memory}) and for
% Chapter~\ref{chap:integration}. Resist adding them here - the first figure is
% the one a reader meets before they know any of the vocabulary.

The first architecture considered was a fully pipelined, layer-per-stage
design, in which every layer of the network is mapped to its own dedicated
hardware stage. Because all stages operate concurrently, such a design
maximizes throughput. Its cost, however, grows with the depth of the network:
every layer requires its own multipliers, buffers and control logic, so a
network with dozens of layers requires dozens of distinct datapaths, each of
which has to be written and verified separately. The area of such a design was
never quantified in detail, but two independent considerations ruled it out:
the resource cost would exceed the target device by a wide margin, and the
development effort would exceed the time available for this work. The approach
was therefore rejected in favor of a program-driven architecture, in which a
central controller sequences a small set of reusable, parameterized compute
engines (\DD{003}).

This choice exploits the structural regularity of MobileNetV2. Although the
network is deep, its layers are instances of only a few operation types: the
standard 3$\times$3 convolution of the first layer (the \emph{stem}),
1$\times$1 pointwise convolutions, 3$\times$3 depthwise convolutions, residual
additions, a final global average pooling, and the linear classifier that
produces the logits. Layers of the same type differ only in their dimensions
and strides. Each type is therefore implemented once, as a parameterized
engine, and reused by every layer of that type. The cost of this reuse is that
only one engine is active at any time; the design thus trades peak throughput
for a resource footprint that fits on the device.

\begin{figure}[t]
  \centering
  \IfFileExists{images/accelerator-design/accel_top.pdf}{%
    \includegraphics[width=\linewidth]{images/accelerator-design/accel_top.pdf}%
  }{%
    \fbox{\parbox{0.8\linewidth}{\centering\vspace{2em}
      \texttt{images/accelerator-design/accel\_top.pdf} not built --- run \texttt{make figures}
    \vspace{2em}}}%
  }
  \caption{The proposed accelerator. A program of \thInstructions{} instructions is
           issued one at a time, and each instruction activates exactly one compute
           engine; because at most one engine is ever live, the logic behind the
           accumulators is shared rather than replicated --- one requantization
           stage and one activation pool, which is read back as the next layer's
           input. The classifier is the one deliberate exception: it reuses the
           pointwise array, but its result is int16 logits rather than an int8
           tensor, so it bypasses both (Section~\ref{sec:sharing}). Weights stream
           from DDR; activations never leave the chip.}
  \label{fig:accel-top}
\end{figure}

Figure~\ref{fig:accel-top} shows the resulting architecture. The network is
expressed as a program of \thInstructions{} instructions, which the program
sequencer issues one at a time. The sequencer decodes each instruction and
presents the resulting configuration to the feeders, which supply the input
data to the selected compute engine in the order that engine expects. The six
operation types are served by five engines --- stem, pointwise, depthwise,
residual and pooling --- because the linear classifier reuses the pointwise
array, as a 1$\times$1 convolution over a single pixel. Each instruction
therefore activates exactly one engine. While that engine is working it holds a
busy signal; when the signal falls and the pipeline has drained, the sequencer
issues the next instruction.

Since at most one engine is active at a time, the logic downstream of the
accumulators does not need to be replicated per engine. The five engines share
a single requantization stage, which rescales the accumulated results back to
int8 and, where the layer requires it, applies the ReLU6 clamp; they also share
a single on-chip activation pool, which stores the output of each layer and is
read back as the input of the next. The classifier is the one deliberate
exception to both: its output is \thClasses{} int16 logits destined for a
register that the processing system reads, not an int8 tensor destined for the
pool, so it ends in a narrow requantization tail of its own.
Section~\ref{sec:sharing} argues why that duplication was accepted rather than
folded into the shared stage. Weights are streamed from off-chip DDR memory
layer by layer, whereas activations remain on-chip for the entire inference.

\section{Design Principle: Engines Perform Only Arithmetic}
\label{sec:engines-philosophy}
% Set on the first RTL day and never broken: the engines take operands and
% produce results; addressing, windowing and sequencing live in the feeders and
% the controller. Two consequences worth stating, because they pay off later:
% every engine is unit-testable against a golden vector with no controller
% present (Chapter~\ref{chap:verification}), and the integration testbenches that build
% windows by hand became the specification of the controller.

From the outset, the compute engines were restricted to arithmetic. An engine
receives operands and produces results; it does not generate addresses, form
convolution windows or decide the order of operations. These tasks are instead
handled by the feeders and the program sequencer.

This separation has two consequences that are exploited later in this thesis.
First, each engine can be verified in isolation: it is tested bit-exactly
against the golden vectors of the int8 software model, without the controller
being present (Chapter~\ref{chap:verification}). Second, the integration
testbenches, which construct the input windows by hand, subsequently served as
an executable specification for the feeders: they define the exact sequence of
operands that the feeders must later generate in hardware.

\section{Parallelizing the Compute Engines}
\label{sec:parallelism}

The first implementation of the engines was sequential: each engine performed
one multiply-accumulate (MAC) operation per clock cycle. A single MobileNetV2
inference costs \thMacsM{} million MAC operations. The total quoted here is the
whole-network figure of Section~\ref{sec:mac-dist}, which includes the global
average pool's accumulations even though they carry no multiply, because a
sequential engine pays a cycle for each of them exactly as it does for a
multiply-accumulate. At one MAC per cycle the baseline therefore needs as many
cycles as it has MACs, a latency of \SI{\thLatencySeq}{\second} at
\MHz{\thClkBoard} --- some \thSpeedupSeq$\times$ slower than the design this
chapter goes on to describe. A sequential design also leaves almost all of the
device's DSP48E2 slices idle, although those slices are the FPGA's principal
resource for parallel arithmetic.

Parallelizing the engines was therefore necessary to reach the performance
target of this work. The same form of parallelism cannot be applied to every
engine, however, because each layer type exposes parallelism along different
loop dimensions. A pointwise convolution reduces over its input channels, and
every input activation is reused by all output channels; parallelism therefore
pays along both dimensions at once, as $T_m = 32$ output channels each
accumulating $T_n = 16$ input channels per cycle --- the $32 \times 16$ array
of Section~\ref{sec:pw-array}. A depthwise convolution, in contrast, convolves
each channel independently with its own filter: no output channel sees more
than one input channel, so there is no cross-channel reduction to parallelize
and the channel count is the only axis available, which makes the depthwise
engine one-dimensional at $T_c = 16$ channels. Whether that engine should be
separate from the pointwise array at all is a further question --- the two
never run concurrently and could in principle share multipliers --- and
Section~\ref{sec:dw-engine} answers it. The following subsections describe each
engine in detail and justify its degree of parallelism.

\subsection{Pointwise Convolution: A \texorpdfstring{$32 \times 16$}{32x16} MAC Array}
\label{sec:pw-array}
% \DD{013}. Tm=32 output lanes x Tn=16 input channels = 512 MAC/cycle, chosen
% against the tiling space of Section~\ref{sec:tiling} and the \thPwMacPct\% MAC share
% of Section~\ref{sec:mac-dist}. Output-stationary; accumulation over ceil(IC/Tn)
% cycles. Broadcast activations, per-lane weights.

Pointwise convolutions account for \thPwMacPct\% of the arithmetic of the
network (Section~\ref{sec:mac-dist}), which is why the pointwise engine was the
first to be parallelized. That share should not be read as a share of the
runtime, however: once the other engines are parallelized as well, this array
accounts for roughly \thPwRuntimePct\% of the modelled inference time, and
Section~\ref{sec:dw-engine} shows how a layer type with a twentieth of the
arithmetic comes to matter nearly as much.

A pointwise convolution is a 1$\times$1 convolution: at each spatial position,
it combines all input channels into each output channel. For an input
activation $x_{i,j,n}$ at spatial position $(i,j)$ and input channel $n$,
output channel $m$ is computed as
\begin{equation}
  y_{i,j,m} = b_m + \sum_{n=0}^{C_{\mathrm{in}}-1} w_{m,n}\, x_{i,j,n},
  \label{eq:pointwise}
\end{equation}
where $w_{m,n}$ is the weight connecting input channel $n$ to output
channel $m$, $b_m$ is the bias of output channel $m$, and $C_{\mathrm{in}}$ is
the number of input channels. Each output element is therefore a dot product
of length $C_{\mathrm{in}}$.

Equation~\eqref{eq:pointwise} can be parallelized along three dimensions: the
input channels, the output channels and the spatial positions. The engine
unrolls the two channel dimensions, $T_m = 32$ output channels and $T_n = 16$
input channels, in the notation of Section~\ref{sec:tiling}, giving
$P = T_m \cdot T_n = 512$ MACs per cycle. The spatial dimension is left
sequential because, with $P$ fixed, unrolling it could only be paid for out of
$T_m$ or $T_n$: pixel parallelism helps only those layers whose output channel
count is smaller than $T_m$, where lanes would otherwise sit idle, and
recovering them was estimated at \thPwPixelGain\% of the total cycle count ---
too little to justify the control logic that a second unrolled dimension
requires (\DD{013}).

With $P$ held at 512, the split between the two channel dimensions is free, and
it matters more than the product does. Evaluating the cycle model
$H W \cdot \lceil OC / T_m \rceil \cdot \lceil IC / T_n \rceil$ over all
\thPwLayers{} pointwise layers of the network gives \thPwUtil\% utilization of
the 512 MAC slots at $32 \times 16$, against \thPwUtilSixteen\% at
$16 \times 32$, \thPwUtilSixtyFour\% at $64 \times 8$ and \thPwUtilEight\% at
$8 \times 64$. The reason lies in the shape of MobileNetV2 itself: the
bottleneck projections are narrow on output, with as few as 16, 24 or 32
output channels, while the expansion layers inside each block are wide on
input. A small $T_m$ with a large $T_n$ matches that asymmetry. The opposite
choice, $T_n = 8$ --- proposed earlier on the grounds that 8 divides every
channel count in the network, so that tail padding would vanish altogether ---
forces $T_m = 64$ with $P$ fixed, and strands most of the lanes on precisely
the layers that are narrowest on output. Weight bandwidth does not discriminate
between the splits, being $T_m \cdot T_n \cdot 8 = 4096$ bit per cycle for all
of them; the remaining tiebreaker is cost, and $32 \times 16$ needs a 128-bit
activation broadcast rather than 256 bits, and an adder tree of depth four
rather than five.

\begin{figure}[t]
  \centering
  % images/accelerator-design/pe_array.pdf is produced from the SVG by `make figures', which
  % needs Inkscape. The guard keeps the manuscript compiling on a machine that
  % does not have it yet - a missing figure should not stop you writing.
  \IfFileExists{images/accelerator-design/pe_array.pdf}{%
    \includegraphics[width=\linewidth]{images/accelerator-design/pe_array.pdf}%
  }{%
    \fbox{\parbox{0.8\linewidth}{\centering\vspace{2em}
      \texttt{images/accelerator-design/pe\_array.pdf} not built --- run \texttt{make figures}
    \vspace{2em}}}%
  }
  \caption{The pointwise MAC array. The $T_m = 32$ lanes share one broadcast of
           $T_n = 16$ activations, while each lane holds its own weights, giving
           $P = T_m \cdot T_n = 512$ MACs per cycle.}
  \label{fig:pe-array}
\end{figure}

As shown in Figure~\ref{fig:pe-array}, the array consists of 32 lanes, one per
output channel. In every cycle, the same 16 input activations are broadcast to
all lanes, and each lane multiplies them by its own 16 weights.

Within each lane, the 16 products are summed by an adder tree, and the
resulting partial sum is added to the lane's accumulator. The accumulator holds
this running sum for $\lceil C_{\mathrm{in}} / T_n \rceil$ cycles, until all
input channels have been processed. Each output element thus stays in place
until it is complete, which makes the engine \emph{output-stationary}. The
array computes only the summation of Equation~\eqref{eq:pointwise}, not the
whole of it: the accumulator is cleared to zero rather than preloaded with the
bias, and $b_m$ is added downstream in the shared output stage, immediately
before requantization. This is the principle of
Section~\ref{sec:engines-philosophy} applied to the one engine where departing
from it would have been easiest and cheapest. The 32 completed output elements
are passed to that stage, and the array proceeds to the next group of output
elements. Neither the lanes nor the array are combinational, and the pipeline
depth they contribute is part of the latency contract of
Section~\ref{sec:latency-contract}. As in the rest of the design, activations
remain on-chip in URAM, whereas weights are streamed from off-chip DDR memory
(Section~\ref{sec:memory}).

The \thPwUtil\% utilization quoted above has a cost worth stating from the
other side: \thPwWastePct\% of the MAC slots that the array issues over a full
inference are multiplications by zero. Whenever a layer's output channel count
is not a multiple of $T_m$, or its input channel count not a multiple of $T_n$,
the final tile is only partly occupied, and it is the feeder rather than the
array that drives zero weights into the unused lanes and taps. Those padded
terms contribute nothing to the sum, so the result remains bit-exact and only
utilization is lost. Keeping the tail in the feeder is what allows the array
itself to stay a fixed $32 \times 16$ block with no notion of a partial tile.

The classifier reuses this array rather than adding a sixth engine, since a
fully connected layer over a feature vector is a 1$\times$1 convolution at a
single spatial position. It is issued as a pointwise instruction over one
pixel, and differs only in its tail: the result is int16 logits rather than
int8 activations, so it bypasses the shared output stage
(Section~\ref{sec:sharing}).

\subsection{Depthwise Convolution: A 16-Channel Engine}
\label{sec:dw-engine}
% \DD{014}, Tc=16. Separate from the pointwise array because the dataflow is
% different, not because of the MAC count. The line buffer reads each input
% pixel ONCE rather than nine times - that is the result worth reporting here.

The depthwise layers are only \thDwMacPct\% of the arithmetic of the network,
and for that reason they were at first treated as a detail. The cycle model
shows this to be the wrong reading. Left at one output element per cycle while
the pointwise array runs at 512 MACs per cycle, the depthwise layers would
account for \thDwShareSeq\% of the inference time: the 512-MAC array would
spend most of its life waiting for a layer type that performs a fifteenth of
the work. Parallelizing the depthwise engine is therefore not the optimization
of a detail --- it is the condition under which the parallelism of the
pointwise array is worth anything at all.

As established in Section~\ref{sec:parallelism}, the channel count is the only
axis available here: each output channel depends on the same-index input
channel alone, so there is no cross-channel reduction to unroll, and widening
the engine means processing more channels at once rather than more terms per
channel. At $T_c = 16$ the depthwise share of the inference falls to
\thDwShare\%, for \thDwDsp{} DSP48E2 slices once the int8 multiplications are
packed two to a slice. The neighbouring choices are what identify this as the
knee of the curve: $T_c = 8$ leaves the depthwise at \thDwShareEight\% of the
runtime, whereas $T_c = 32$ reduces it only to \thDwShareThirtyTwo\% in
exchange for twice the slices (\DD{014}).

That the depthwise engine is separate from the pointwise array is a second
decision, and it does not follow from the first. Because the network alternates
pointwise, depthwise and pointwise layers, the two engines never run
concurrently, and they could in principle share one pool of multipliers. They
do not, because the two dataflows have nothing in common: a broadcast
activation bus feeding a cross-channel adder tree on one side, sixteen
independent nine-tap dot products on the other. Merging them would buy
multipliers at the price of multiplexing the one resource that is not the
binding constraint of this design --- the finished accelerator uses
\thDspPctAccelAxi\% of the device's DSP slices --- and would pay for them in
control complexity and timing margin on the critical engine.

\begin{figure}[t]
  \centering
  % images/accelerator-design/dw_array.pdf is produced from the SVG by `make
  % figures'. Same guard as the other two figures: a missing PDF should not
  % stop the manuscript compiling on a machine without Inkscape.
  \IfFileExists{images/accelerator-design/dw_array.pdf}{%
    \includegraphics[width=\linewidth]{images/accelerator-design/dw_array.pdf}%
  }{%
    \fbox{\parbox{0.8\linewidth}{\centering\vspace{2em}
      \texttt{images/accelerator-design/dw\_array.pdf} not built --- run \texttt{make figures}
    \vspace{2em}}}%
  }
  \caption{The depthwise array. Sixteen independent lanes, one per channel,
           each nine multipliers and an adder tree; there is no accumulator,
           because a depthwise dot product is exactly nine taps and completes
           in one pass, and no connection of any kind between lanes. The line
           buffer delivers the window tap-major, so lane $c$ gathers byte $c$
           of every tap at a stride of $T_c$ bytes, while the channel-major
           weight banks pass straight through.}
  \label{fig:dw-array}
\end{figure}

Figure~\ref{fig:dw-array} shows the resulting engine: sixteen independent
lanes, one per channel, each an instance of the same 3$\times$3 convolution
unit, in which nine multipliers produce nine products that a balanced adder
tree reduces to one sum. Two things distinguish a depthwise lane from the
pointwise lane of Section~\ref{sec:pw-array}. First, it has no accumulator: a
depthwise dot product is exactly $K \cdot K = 9$ taps and completes in a single
pass, where a pointwise lane streams $\lceil C_{\mathrm{in}} / T_n \rceil$
tiles through its accumulator. Second, the lanes never interact --- no adder
tree spans them and no activation is broadcast between them --- which makes the
engine, structurally, a pointwise array turned inside out. Altogether it holds
$T_c \cdot K \cdot K = 144$ multipliers. The adder tree is built for sixteen
leaves rather than nine, since nine is not a power of two; the seven unused
taps are padded with zeros on the buses rather than special-cased in the logic,
and synthesis prunes the constant products and the additions that consume them.

Almost all of the array's own logic is a transposition between two bus orders,
and it is worth setting them out. The line buffer delivers the window
tap-major and channel-minor, as $\mathrm{win}[(k_y K + k_x) T_c + c]$, which is
the natural order for it to produce: one window position is $T_c$ channels
wide, because all sixteen channels advance together. The weight buffer delivers
the opposite order, channel-major and tap-minor, which is equally natural for
it, because bank $c$ \emph{is} channel $c$ and holds that channel's nine
weights contiguously. Each lane, however, needs the nine taps of one channel.
The weights therefore arrive already contiguous per lane and pass straight
through, while the window has to be gathered with a stride of $T_c$ bytes.
Getting that stride wrong would mix channels while still producing entirely
plausible sums, with no width mismatch and no elaboration error to betray it,
which is what the array's testbench is written to catch
(Chapter~\ref{chap:verification}).

Consistently with Section~\ref{sec:engines-philosophy}, none of the window
construction belongs to the array. The line buffer and the feeder generate the
window, and zero padding at the image borders is the feeder's responsibility
as well: it drives a zero into any tap that falls outside the image, which the
lane multiplies by its weight and contributes nothing, exactly as the
reference C model skips that tap. The line buffer is also what makes the
dataflow efficient, since it retains the rows the window needs and therefore
reads each input pixel once rather than nine times.

One consequence of generating windows this way is easy to overlook, and was
overlooked in the first version of the cycle model: the cost is driven by the
input pixels, not the output pixels. A sliding-window generator must consume
every input pixel in order to advance the window, including at stride two,
where only every second window yields an output. \thDwStrideTwo{} of the
\thDwLayers{} depthwise layers have stride two, and their input grids are four
times the size of their output grids, so counting output pixels understated
the depthwise cost by a factor of \thDwInputFactor{}. The figures quoted above
are from the corrected model; the choice of $T_c = 16$ survived the correction
unchanged.

The accumulator width, finally, is a provable bound here rather than a
data-dependent precondition. A depthwise dot product is always exactly nine
taps, so $\lvert \mathrm{acc} \rvert \le 9 \cdot 128 \cdot 127 = \num{146304}$
for any int8 input whatsoever, which needs 19 bits and sits well inside the
26-bit accumulator. Section~\ref{sec:widths} returns to this, and contrasts it
with the pointwise side, where the same width is arrived at by a different
argument.

\subsection{Stem Convolution: An 8-Channel Engine}
\label{sec:stem-engine}
% Runs once per inference, on 3 input channels. Sized accordingly; justify the
% asymmetry with \thStemMacPct\%.

The stem is the first layer of the network, \file{features.0.0}: a standard
3$\times$3 convolution at stride 2 with padding 1, taking the
$224 \times 224 \times 3$ image to a $112 \times 112 \times 32$ feature map,
with a ReLU6 activation. It is the only standard convolution in MobileNetV2 ---
every other convolutional layer is either pointwise or depthwise --- and it is
also the only operation whose input is the image rather than the activation
pool. Because it consumes the image directly, its input channel count is fixed
at three, the colour channels, and each output element is therefore a dot
product of exactly $C_{\mathrm{in}} \cdot K \cdot K = 27$ taps. One lane of the
engine holds 27 multipliers, which makes the degree of parallelism a genuine
area decision rather than a free one.

That decision has to answer an apparent asymmetry: the pointwise array is
given 512 multiply-accumulate operations per cycle and the stem only eight
lanes. The stem accounts for \thStemMacPct\% of the arithmetic of the network
and executes once per inference, so it does not warrant a comparable share of
the device. Left at one output element per cycle, however, it is not
negligible either --- it would be \thStemShareSeq\% of the modelled inference
time, which is far more than its share of the arithmetic, for the same reason
the depthwise layers were in Section~\ref{sec:dw-engine}.

What fixes the number at eight is the rate of the input stream rather than the
arithmetic. The line-buffer front end streams $225 \times 225$ input
pixel-groups, the image plus its virtual border, and emits $112 \times 112$
windows at stride 2, so there are exactly $50\,625 / 12\,544 = 4$ input cycles
per window. Each window needs $\lceil OC / T_s \rceil$ cycles of compute, so
$T_s \ge 8$ is the point at which the compute stops exceeding the input
stream; below it the line buffer idles while the array sweeps, and above it
the array idles waiting for windows. The measured sweep confirms the knee:
$T_s = 4$ leaves the stem at \thStemShareFour\% of the runtime, $T_s = 8$
brings it to \thStemShare\%, and $T_s = 16$ reaches only \thStemShareSixteen\%
for twice the multipliers (\DD{016}). It is worth being precise about what
widening the engine bought, because the arithmetic invites the wrong answer:
going from one lane to eight made the stem roughly
\thStemSpeedup$\times$ faster, not eight times. A linear speedup would assume
the compute remained the limit throughout, and the entire point of choosing
eight is that at eight it no longer is.

\begin{figure}[t]
  \centering
  \IfFileExists{images/accelerator-design/stem_array.pdf}{%
    \includegraphics[width=\linewidth]{images/accelerator-design/stem_array.pdf}%
  }{%
    \fbox{\parbox{0.8\linewidth}{\centering\vspace{2em}
      \texttt{images/accelerator-design/stem\_array.pdf} not built --- run \texttt{make figures}
    \vspace{2em}}}%
  }
  \caption{The stem engine. Eight lanes each compute one output channel from the
           \emph{same} 27-tap window of $3 \times 3 \times 3$: a standard
           convolution sums over the input channels as well as the window, so the
           activations are broadcast as in the pointwise array of
           Figure~\ref{fig:pe-array} rather than split per lane as in the depthwise
           array of Figure~\ref{fig:dw-array}. The window is transposed into
           channel-major order once, before the broadcast, not once per lane. With
           $OC = 32$ and $T_s = 8$ each window is swept in four oc-tiles, one per
           cycle, which is exactly the rate at which the line buffer produces the
           next window; the four groups of eight are assembled into one 32-channel
           entry. The stem is the only operation that reads the image buffer rather
           than the activation pool.}
  \label{fig:stem-array}
\end{figure}

Figure~\ref{fig:stem-array} shows the engine. Its eight lanes are instances of
the same standard-convolution unit, each producing one output channel, and
they are arranged like the pointwise array rather than like the depthwise one.
A standard convolution sums over the input channels as well as over the
window, so every output channel is computed from the same 27 activations and
the lanes differ only in their weights: the window is broadcast, exactly as the
activation tile is broadcast in Figure~\ref{fig:pe-array}, and not split per
lane as in Figure~\ref{fig:dw-array}. The same transposition problem described
in Section~\ref{sec:dw-engine} arises here, since the line buffer again emits
tap-major and the convolution unit wants channel-major order, but because the
window is shared it is transposed once, ahead of the broadcast, instead of
once per lane. The weights need no transposition at all: the manifest stores
them as $(OC, IC, K_h, K_w)$ row-major, which is already the order the unit
expects. As in the depthwise lane there is no accumulator, because 27 taps is
the whole dot product and it completes in one pass. In total the engine holds
$T_s \cdot C_{\mathrm{in}} \cdot K \cdot K = 216$ multipliers, or
\thStemDsp{} DSP48E2 slices with the int8 products packed two to a slice.

Since $OC = 32$ and $T_s = 8$, a window is swept in four oc-tiles, one per
cycle, the lanes computing output channels $8t$ to $8t + 7$ on tile $t$. The
four results of eight channels are assembled into a single 32-channel entry
before it is written to the activation pool, so the engine emits one complete output pixel per window.
This is also where the balance of the previous paragraph becomes visible in
the control logic: the windows do not arrive evenly, but every two cycles
inside an odd row and not at all on even rows, so the feeder stalls the input
stream while each window is swept. The lane array itself knows nothing of
this. Consistently with Section~\ref{sec:engines-philosophy}, the counters,
the window generation, the border padding and the stall all belong to the
feeder; the array receives 27 activations and eight sets of weights, and
returns eight sums.

One consequence of the 27-tap dot product belongs to Chapter~\ref{chap:timing}
rather than here, but is worth signposting: this unit, not the much larger
pointwise array, was the critical path of the design from the first synthesis
onward, because 27 taps in an unregistered cascade is a far deeper chain of
logic than anything in the 512-MAC array. Pipelining it was the single largest
improvement in the timing work.

\subsection{Requantization Lanes}
\label{sec:requant-lanes}
% The shared out_stage. \DD{016}: sharing it took the predicted DSP count from
% ~634 to 442. Implements Section~\ref{sec:requant-theory} exactly.
% NOTE: the DD number above is wrong. The sharing argument is \DD{015} and
% out_stage.v's header; DD-016 is the stem's TS. Left as written so the error
% is visible; the text below cites \DD{015}.

Every engine described so far produces accumulators, not activations. A
pointwise lane accumulates into 26 bits, and the next layer consumes
int8, so each result must be returned to the int8 grid before it can be
stored. That conversion is requantization, and it is the only arithmetic in
the accelerator that is neither a multiplication nor an addition of the
operands themselves.

The quantization scheme is int8 with per-channel weight scales and per-tensor
activation scales, and batch normalization folded into the weights and biases
offline (\DD{005}). Two distinct models are involved, and the distinction
matters for what each one is evidence of. The \emph{fake-quantized} model runs
in floating point and merely rounds values onto the int8 grid at each step; it
was used to establish that the scheme does not cost accuracy, but its
floating-point rounding is not reproducible in hardware to the last bit, so it
is not a specification. The \emph{integer-only} model, written separately,
performs every operation on integers exactly as the RTL does, and it is this
second model that defines what the hardware must compute. The accelerator is
verified against the integer-only model alone; the fake-quantized model never
appears downstream of accuracy measurement.

Let $S_x$, $S_w$ and $S_y$ be the scales of the input activations, the weights
of one output channel and the output activations, all three being the positive
real number by which an integer code is multiplied to recover the value it
represents. An accumulator $a$ holds the integer sum of products of int8
codes, so the value it stands for is $S_x S_w a$, whereas the output code
must be expressed in units of $S_y$. The conversion is therefore a
multiplication by the real ratio $M = S_x S_w / S_y$, which is per output
channel because the weight scales are. Since the hardware has no floating
point, $M$ is represented exactly as a fixed-point pair, and the operation
performed on each element is
\begin{equation}
  y \;=\; \mathrm{clamp}\!\left(
          \left\lfloor \frac{a \cdot M_0 + 2^{\,s-1}}{2^{\,s}} \right\rfloor
          \right),
  \qquad M \;\approx\; M_0 \cdot 2^{-s},
  \label{eq:requant}
\end{equation}
where $M_0$ is a signed integer multiplier and $s$ the right-shift amount.
The exporter normalizes the pair so that $2^{30} \le M_0 < 2^{31}$, which fixes
$s \ge 31$ for every layer of this network, because every real ratio $M$ in it
is smaller than one.

Three details of Equation~\eqref{eq:requant} are load-bearing, in the sense
that changing any of them changes results. First, the floor is an
\emph{arithmetic} right shift, and the constant $2^{s-1}$ added before it is
half an output LSB, so the rule is round-half-up toward $+\infty$: a negative
tie rounds up, $-2.5$ giving $-2$ rather than $-3$. Second, the product
$a \cdot M_0$ is formed at full width and never truncated. Third, the clamp
depends on the activation function of the layer. A layer with ReLU6 clamps to
$[0, q_6]$, where $q_6 = \min(127, \mathrm{round}(6 / S_y))$ is the integer
code of the value 6 at the output scale; because $q_6 \le 127$, this clamp
subsumes the int8 saturation and no separate saturation is applied. A layer
without one saturates to $[-128, 127]$. The classifier is the same arithmetic
with a wider saturation, to $[-32768, 32767]$, since its output is int16
logits rather than int8 activations (\DD{011}).

\begin{figure}[t]
  \centering
  \IfFileExists{images/accelerator-design/requantize.pdf}{%
    \includegraphics[width=\linewidth]{images/accelerator-design/requantize.pdf}%
  }{%
    \fbox{\parbox{0.8\linewidth}{\centering\vspace{2em}
      \texttt{images/accelerator-design/requantize.pdf} not built --- run \texttt{make figures}
    \vspace{2em}}}%
  }
  \caption{One requantization lane, as four register stages. Each parameter
           enters at the stage that consumes it rather than at the input, so the
           window in which a stale parameter can meet live data is one cycle
           wide regardless of the pipeline depth. Only the clamp narrows the
           datapath.}
  \label{fig:requantize}
\end{figure}

Figure~\ref{fig:requantize} shows how Equation~\eqref{eq:requant} is realized.
The bias addition and the requantization form four register stages: the biased
accumulator, the multiplication by $M_0$, the rounding and shift, and the
clamp. The intermediate is carried at
$\mathit{ACC\_W} + \mathit{M0\_W} + 1 = 26 + 32 + 1 = 59$ bits, and the clamp
in the fourth stage is the only point at which the datapath narrows. Every
operand in the chain is declared signed; a single unsigned operand would make
the whole expression unsigned and silently break both the arithmetic shift and
the comparisons.

The stage at which each parameter enters is itself a decision. The bias, $M_0$
and the shift are per output channel and arrive from the parameter buffer,
while the activation kind and $q_6$ are per layer and come from the
instruction word; originally all of them were presented at the input, which
was correct only because the sequencer happens not to fetch the next
instruction while an operation is running. With one pipeline stage that
coincidence was harmless. With four it protects an arithmetic unit four levels
below the place that guarantees it, so each parameter is now delayed to the
stage that uses it: $M_0$ to the second, the shift to the third, and the
activation parameters to the fourth. The guarantee became structural rather
than contractual, at a cost of a few registers.

How many such lanes to build was the one sizing decision here. The pointwise
array finishes $T_m = 32$ accumulators every $\lceil IC/T_n \rceil$ cycles, and
with $R$ lanes the drain takes $\lceil T_m / R \rceil$ cycles, so an output
tile costs $\max(\lceil IC/T_n \rceil, \lceil T_m/R \rceil)$ cycles: too few
lanes and the array stalls waiting to unload. Modelled over all
\thPwLayers{} pointwise layers, $R = 1$ costs \thRqPenaltyOne\% more cycles
than an unstalled design, $R = 8$ costs \thRqPenaltyEight\% and stalls
\thRqStalledEight{} layers, and $R = 16$ costs only \thRqPenaltySixteen\% and
stalls \thRqStalledSixteen. The design uses $R = T_m = 32$, one lane per
accumulator, for \thRqDsp{} DSP48E2 slices (\DD{015}).

$R = 16$ is the knee of the cycle curve and was the obvious alternative, but it
was rejected for a reason that the cycle counts do not show. At $R = 32$ all
32 results appear in the same cycle and already constitute the 256-bit word
that the activation buffer stores. At any smaller $R$ the output stage needs a
shadow register to hold the accumulators while the array moves on, a drain
counter, a stall path back into the address generator, and logic to assemble
the stored word in pieces --- an entire control mechanism, and one that
$R = 16$ would exercise on exactly one layer of the \thPwLayers{}. The 32
additional slices were judged the cheaper side of that trade.

A second decision made this block shared rather than replicated. Because the
sequencer runs one operation at a time, no two engines ever requantize
simultaneously, so the bank and its parameter memory were factored out of the
feeders into a single stage used by all of them. Before that, the pointwise
path carried a 32-lane bank and the depthwise path a 16-lane one, and the stem,
residual and pooling paths were each about to add their own. Sharing was
predicted to reduce the DSP budget from roughly \thDspPredUnshared{} to
\thDspPredShared{} and, as importantly, left one block to verify instead of
five. Both figures are predictions made before the timing work;
Chapter~\ref{chap:evaluation} reports the measured utilization of the finished
design. The classifier is the single exception to the sharing, for the reason
given in Section~\ref{sec:sharing}.

The lane is verified directly against the integer-only model rather than
against a tolerance. Its unit testbench replays the same directed vectors as
the C model's own test of its requantization routine, so a pass establishes
that the RTL agrees with the reference on every case the reference is itself
validated against, including the rounding of negative ties and both clamps.
The bank and the whole-network run extend this to every element of every
layer, as Chapter~\ref{chap:verification} describes.

One divergence between the two is deliberate and is worth recording, since a
later reader will otherwise read it as a defect. The C model narrows its
intermediate to 32 bits before clamping; the RTL clamps the full 59-bit value.
The two therefore disagree whenever the shifted result exceeds $2^{31}$ in
magnitude, where the C model wraps and the hardware saturates. That region is
unreachable, and by a bound rather than by observation: the normalization that
gives $2^{30} \le M_0 < 2^{31}$ also forces $s \ge 31$, so
\begin{equation}
  \left| a \cdot M_0 \cdot 2^{-s} \right| \;\le\; |a| \cdot \frac{M_0}{2^{s}}
  \;<\; |a| \;\le\; 2^{25},
  \label{eq:requant-bound}
\end{equation}
which is six bits clear of $2^{31}$ for any input data whatsoever.
Equation~\eqref{eq:requant-bound} depends only on the exported multipliers
being normalized, not on any measurement, which is the distinction
Section~\ref{sec:acc-bound} draws for the accumulator width. The RTL is the
more conservative of the two and nothing reads a wider value from the C model,
so the divergence was documented and kept rather than removed, and the
simulation model carries a check that fires if it is ever reached (\DD{024}).

\subsection{Datapath Widths}
\label{sec:widths}
% Every width, and for each one whether it is a BOUND or a MEASUREMENT. After
% \DD{019} this table is one of the more defensible things in the thesis -
% make the distinction a column, not a footnote.

In an integer accelerator every signal has a width, and each width is justified
in one of two ways: by a bound, which holds for every input the format admits,
or by an observation, which holds for the inputs that happened to be tried.
The distinction is not academic. One width in this design was once chosen from
an observation, and the section below describes what that cost. Table~%
\ref{tab:widths} lists every width in the datapath together with what fixes it
and which kind of justification it rests on. Storage geometry --- buffer
depths, banking and address widths --- is a separate matter and is treated in
Section~\ref{sec:buffers}.

\begin{table}[t]
  \centering
  \caption{The datapath widths, and what justifies each. A \emph{bound} holds
           for every input the format admits; an \emph{observation} holds for
           the inputs that were tried. Only the logit scale now rests on the
           latter.}
  \label{tab:widths}
  \small
  \begin{tabular}{llll}
    \toprule
    Signal & Width & Fixed by & Justification \\
    \midrule
    Activations, weights   & 8  & the quantization scheme (\DD{005})      & construction \\
    Tap product            & 16 & exact int8 $\times$ int8                 & exact \\
    Accumulator            & 26 & worst-case $|acc + b|$, any input (\DD{019}) & bound \\
    Bias                   & 32 & int32, as exported                       & construction \\
    Requant.\ multiplier   & 32 & normalized to $[2^{30}, 2^{31})$         & bound \\
    Requant.\ shift        & 6  & largest exported shift is 48             & bound \\
    Requant.\ intermediate & 59 & $\mathit{ACC\_W} + \mathit{M0\_W} + 1$   & exact \\
    Activation output      & 8  & the quantization scheme                  & construction \\
    Logit output           & 16 & $S_{\mathrm{logit}}$, calibrated (\DD{011}) & observation \\
    \bottomrule
  \end{tabular}
\end{table}

The width of the activations and the weights is 8 bits by the choice of
quantization scheme (\DD{005}), and the width of the output is the same,
because the output of one layer is the input of the next. The product of two
int8 operands is exact in 16 bits, since the extreme case $-128 \times -128$
is $16\,384$, inside the signed 16-bit range; the engines therefore register
each tap product at that width before the adder tree, and no rounding occurs
anywhere between the multiplier and the accumulator. The bias is int32 as the
exporter emits it, and the widths internal to requantization --- the 32-bit
multiplier, the 6-bit shift and the 59-bit intermediate --- are those of
Equation~\eqref{eq:requant}, discussed in Section~\ref{sec:requant-lanes}. The
6-bit shift field deserves one remark: it is a bound rather than a convenience,
because the exporter's normalization caps the shift well below the 63 that six
bits can encode, the largest value in the current network being 48.

The accumulator is the width this design got wrong, and it is worth setting out
why, because the failure was methodological rather than arithmetic. The width
was originally 21 bits, chosen from a measurement: the largest
$|acc + b|$ observed while running one image through one checkpoint was
\thAccMeasured, which needs \thAccMeasuredBits{} signed bits and left what
looked like comfortable headroom. The model was then retrained. Re-exporting
the weights moved one channel of \file{features.3.conv.0.0} to a bias of
\thBadBias, which does not fit a 21-bit signed accumulator \emph{on its own},
for any input whatsoever. The hardware wrapped it negative, the ReLU6 clamp
turned the wrapped value into zero, and \thBadElements{} of the
\thElements{} elements of the network were wrong. Nothing in the RTL detects
accumulator overflow, so the only symptom was the golden-vector comparison
failing; the design had no mechanism that could have caught it earlier.

The correction was to size the accumulator from a bound. For one output
channel $oc$ of a layer, the accumulator before requantization satisfies
\begin{equation}
  \left| acc_{oc} + b_{oc} \right| \;\le\;
  \left( \sum_{\mathrm{taps}} \left| w_{oc,\mathrm{tap}} \right| \right)
  \cdot \max |a| \;+\; \left| b_{oc} \right| ,
  \label{eq:acc-bound-widths}
\end{equation}
where $w_{oc,\mathrm{tap}}$ are the int8 weights of that channel, $b_{oc}$ its
int32 bias, and $\max |a|$ the largest magnitude an input activation can take.
The last quantity is what makes Equation~\eqref{eq:acc-bound-widths} a bound
rather than a statistic: an activation entering a layer is either the output of
a ReLU6, and so lies in $[0, q_6]$, or a plain int8 in $[-128, 127]$. Both
limits are properties of the number format, not of the data, so no image ---
including one chosen adversarially --- can exceed them.

Evaluated over this network, Equation~\eqref{eq:acc-bound-widths} peaks at
\thAccBound{} in the classifier, which requires \thAccBoundBits{} signed bits.
Two things about that figure are worth noting. The first is that it is more
than four times the measured peak, which is the usual relationship between a
bound and an observation. The second is more pointed: the bound and the
measurement identify \emph{different layers}. The measured worst case was
\file{features.3.conv.0.0}, an outlier bias, whereas the provable worst case is
the classifier, which simply has the longest dot product in the network.
Sizing from the measurement would have been unsound even if the measured number
had been larger, because it was taken at the wrong place. A weaker bound that
discards the weights entirely and uses only the tap count and the format,
$\mathrm{taps} \cdot 128 \cdot 128$, gives \thAccGeom{} for the same layer and
is also instructive: it is reached by no real weight set, yet the chosen width
accommodates even that.

The design uses 26 bits, whose range of $\pm\thAccRange$ leaves
\thAccMargin$\times$ margin over the bound. The choice of 26 rather than the
\thAccBoundBits{} the bound strictly requires is a hardware argument. The
requantization multiply is $\mathit{ACC\_W} \times \mathit{M0\_W}$, and the
DSP48E2 multiplier is $27 \times 18$ bits. While the accumulator is at most 27
bits wide, one operand fits a port whole and the 32-bit multiplier splits
across two slices, costing two slices per lane; at 28 bits both operands must
be split and the cost doubles to four, which over the $R = 32$ lanes of
Section~\ref{sec:requant-lanes} would double the \thRqDsp{} slices spent there,
bought for no correctness. Everything up to 27 bits is therefore free and 28 is
expensive, so the width was set at the top of the free range rather than at
the bottom of the safe one. Taking the minimum of \thAccBoundBits{} bits would
have left only \thAccMarginNarrow$\times$ margin for no saving, and 32 bits
would have bought nothing the bound does not already guarantee while paying
the doubled multiplier (\DD{019}).

The cost of the extra five bits was measured with two synthesis runs at the
same timing constraint, so that the width was the only variable. The
accumulator widening added \thAccLutDelta{} CLB LUTs, or \thAccLutPct\%, spread
across the feeder multiplexer, the delay lines and the wider adders, and left
the DSP, block-RAM and URAM counts unchanged. The DSP count is the figure the
decision turned on, and it did not move, which is the whole content of the
27-bit argument above. Timing did not degrade either, since the critical path
at the time lay in the tap cascade of Section~\ref{sec:stem-engine} and not in
the accumulator.

Two alternatives were considered and rejected. Re-quantizing the model so that
the biases fit the narrower accumulator would have changed the numerics for a
reason unrelated to accuracy, and would have left the identical problem waiting
for the next retrain. Saturating the accumulator instead of letting it wrap
would have converted a silent corruption into a visible one, which is an
improvement, but a design that saturates is still a design that has run out of
bits. Testing more images was rejected on principle: more observations raise
confidence without establishing a bound, and in this instance they would have
found nothing at all, because the overflow was contained in a bias and was
therefore independent of the input.

What replaced the measurement is a check rather than a number. The exporter
now evaluates Equation~\eqref{eq:acc-bound-widths} over every output channel of
every layer, reads the accumulator width out of the RTL source so that the two
cannot drift apart, and exits with an error if the network's worst case does
not fit. Re-running that check against the current export confirms the figures
quoted above, and also shows how far the original width was from safe:
at 21 bits, \thAccUnsafeLayers{} layers of the network exceed the bound. Only
one of them produced an observable error on the test image, which is precisely
why the observation was misleading.

One width is still justified by observation, and the table records it as such.
The classifier emits int16 logits on a shared scale $S_{\mathrm{logit}}$,
calibrated from the largest logit magnitude the float model produced over a
few batches, with a factor of 1.2 in hand (\DD{011}). A shared scale is what
lets the predicted class be taken by comparing the \thClasses{} logits
directly, with no rescaling. The exposure is bounded in practice because the
requantization saturates rather than wraps, as Section~\ref{sec:requant-lanes}
describes: a logit beyond the calibrated range is clipped to the end of the
int16 interval, which costs resolution at the extremes and could in principle
make two saturated classes compare equal, but cannot corrupt the result the way
a wrapped accumulator does. The margin is wide --- \thClasses{} classes on a
16-bit scale --- so the risk is small, but it is an observation, and
Table~\ref{tab:widths} says so.

\section{Memory Organization}
\label{sec:memory}

\subsection{On-Chip Activations and Off-Chip Weights}
\label{sec:mem-split}
% Activations never leave the chip; weights stream from DDR. Justify from
% Section~\ref{sec:roofline-analysis}: \thWeightMB{} MB of weights against activations
% that fit in URAM.

\subsection{Activation Pool and Static Allocation}
\label{sec:act-pool}
% One pool, statically allocated by the compiler. No allocator at run time, so
% no fragmentation and no dynamic failure mode.

\subsection{Weight, Parameter and Image Buffers}
\label{sec:buffers}
% Sizes, banking, and the address-width headroom noted in Section~\ref{sec:risks}:
% features.18.0 uses 800 of 1024 entries, and the classifier is small only
% because this model has \thClasses{} classes.

\section{Control Unit and Instruction Set}
\label{sec:control}

\subsection{Program Sequencer}
\label{sec:sequencer}
% Sequential issue, one op at a time, \thInstructions{} instructions. State the
% ceiling plainly: 64 of 64 used, zero spare, and it fails at build time rather
% than corrupting execution.

\subsection{Instruction Format and Program Compiler}
\label{sec:isa-format}
% scripts/gen_program.py compiles manifest.json to the program. Full encoding
% in Chapter~\ref{appendix:isa}; here, only the shape and the six opcodes.

\subsection{Feeders and Handshaking}
\label{sec:feeders}
% Five feeders, one per dataflow. The addr_gen contract - shape latched at
% start so the caller may prefetch - and the C-3 finding: no feeder honoured
% it, masked only because top_seq does not refetch until !op_busy. Worth
% telling, because implementing the advertised prefetch would have broken all
% five silently, as wrong addresses rather than as an error.

\subsection{Resource Sharing}
\label{sec:sharing}
% \DD{016}, \DD{017}. out_stage shared across engines; res_feeder needs no
% multiplier; gap_feeder needs no DSP; the classifier needs no feeder at all.
% \DD{017} - two busy bugs that appeared only when sharing became real - is
% evidence for Section~\ref{sec:three-levels} and should be cross-referenced there.

\section{Latency as a Contract}
\label{sec:latency-contract}
% SOURCE: docs/latency_map.md. Pipeline depths are restated in five feeders, so
% a kernel cannot change in isolation. This is the structural fact that made
% the timing work in Chapter~\ref{chap:timing} expensive, and A-2 (dw_array exporting
% valid/valid_out so latency is a wire rather than a constant) is the fix
% pattern. Put the latency table here and refer to it from Chapter~\ref{chap:timing}.
